\chapter{Acanthosis Nigricans}

\section*{Definition}
Acanthosis nigricans (AN) is a skin disorder characterized by symmetric, thickened, darker, velvety plaques in body folds such as the neck, axillae, and groin. It is commonly associated with obesity or insulin resistance, but can also be familial, related to an endocrine disorder or medicine, or, rarely, associated with internal malignancy. Similar darkening can occur in other conditions, so the appearance alone does not establish the cause.

\section*{When to See a Doctor}

\begin{itemize}
 \item Sudden onset, rapid progression, extensive disease, severe itching, or involvement of the mouth, lips, palms, soles, or other unusual sites, particularly without an obvious metabolic or endocrine cause
 \item Unintended weight loss, loss of appetite, enlarged lymph nodes, persistent gastrointestinal symptoms, or other symptoms suggesting an internal illness
 \item New or changing dark, thickened skin should be assessed rather than self-diagnosed; a skin biopsy may be needed when the appearance is atypical
\end{itemize}

\section*{How It Develops}
In insulin resistance, higher insulin levels may stimulate insulin and insulin-like growth factor signaling in skin cells. This can promote epidermal thickening and papillomatosis, producing the characteristic velvety texture and darker appearance. In malignant AN, tumor-related growth signals can produce similar changes. The underlying mechanism is complex and cannot by itself identify whether a lesion is benign or malignant.

\section*{Causes}
\begin{itemize}
 \item \textbf{Metabolic:} Obesity, insulin resistance, prediabetes, type 2 diabetes, metabolic syndrome, and polycystic ovary syndrome (PCOS)
 \item \textbf{Genetic or familial:} Familial AN, insulin-receptor disorders, and syndromes such as Alstr\"om syndrome or Down syndrome
 \item \textbf{Malignant:} Rare paraneoplastic AN, most often associated with gastrointestinal adenocarcinoma but also reported with other cancers, including liver, lung, ovarian, and hematologic malignancies
 \item \textbf{Medication-induced:} Some systemic glucocorticoids, niacin, growth hormone, and hormonal medicines, including some oral contraceptives
 \item \textbf{Endocrine:} PCOS, Cushing's syndrome, acromegaly, thyroid disease, or adrenal disease
\end{itemize}

\section*{Related Symptoms}
\begin{itemize}
 \item Obesity or increased waist circumference, skin tags (acrochordons), high blood glucose, hirsutism or irregular periods in PCOS, and other signs of insulin resistance
\end{itemize}
\textbf{Important:} Treatment should focus on an identified underlying condition. Improvement may be gradual, and AN can persist even after blood glucose or weight improves. Sudden widespread or atypical AN warrants medical assessment; it does not automatically mean that cancer is present.

\section*{Medical Evaluation}
\begin{itemize}
 \item A clinician usually diagnoses AN by examining the skin; a biopsy is reserved for uncertain or atypical cases.
 \item Blood pressure, body measurements, fasting glucose or HbA1c, and a lipid profile may be used to assess metabolic risk. Fasting insulin is not routinely required to diagnose insulin resistance or diabetes.
 \item Evaluate for PCOS when symptoms such as irregular periods, excess hair growth, or acne are present; hormone tests and pelvic ultrasound are selected according to the clinical assessment rather than required for everyone.
 \item If rapid, extensive, atypical AN or associated symptoms raise concern for malignancy, evaluation is individualized by age, symptoms, examination, and risk factors. This may include targeted imaging, endoscopy, or other tests; routine CT, upper endoscopy, and colonoscopy are not indicated for every person with AN.
\end{itemize}

\section*{Treatment Options}
\begin{itemize}
 \item Treat the underlying condition. Healthy eating, physical activity, and medically appropriate weight management may improve insulin resistance; metformin is used for selected conditions such as diabetes or PCOS, not routinely just to treat the skin appearance.
 \item Topical treatments such as tretinoin, ammonium lactate, urea, or other keratolytic medicines may modestly improve texture or color and should be selected and monitored by a clinician because they can irritate skin.
 \item Laser or other procedures may be considered by a dermatologist for persistent cosmetic concerns, but evidence and benefit vary and pigment changes or scarring are possible. Cryotherapy is not routine treatment for AN.
 \item Drug-induced AN may improve after the responsible medicine is changed under medical supervision. Malignant AN is managed by treating the underlying cancer.
\end{itemize}

\section*{Self-Care and Home Management}
\begin{itemize}
 \item If appropriate for your health, regular physical activity and sustainable nutrition changes can improve metabolic health; weight loss is not necessary or appropriate for everyone.
 \item A plain moisturizer may reduce dryness and friction. Products containing ammonium lactate or urea can irritate some skin and should be used as directed.
 \item Avoid repeated friction, harsh scrubbing, bleaching products, or irritation of affected skin; these do not remove AN and may worsen discoloration.
 \item Do not start metformin, bleaching creams, or prescription retinoids without medical advice.
\end{itemize}

\section*{References}
\begin{thebibliography}{99}
\bibitem{acanthosis_nigricans:ref1} American Academy of Dermatology. Acanthosis Nigricans: Who Gets It and Causes; Treatment; Self-Care. Current online resources.
\bibitem{acanthosis_nigricans:ref2} MedlinePlus. Acanthosis Nigricans. U.S. National Library of Medicine. Reviewed 2026.
\bibitem{acanthosis_nigricans:ref3} DermNet. Acanthosis Nigricans: Causes, Diagnosis, and Treatment. Updated 2024.
\bibitem{acanthosis_nigricans:ref4} James WD, Elston DM, Treat JR, Rosenbach MA, Neuhaus IM. \textit{Andrews\' Diseases of the Skin: Clinical Dermatology}. 14th ed. Elsevier, 2026.
\end{thebibliography}
